\documentclass{situationist-preprint}

\title{Quine's Proposal Reconsidered}
\msid{EPI-26-b}
\msfamily{EPI-26}
\msdiscipline{Epistemology}
\mstype{Companion manuscript}
\mskeywords{naturalized epistemology; Quine; Kim; normativity; engineering; bridge principles}

\begin{document}
\maketitle

\begin{abstract}
Quine's proposal that epistemology become a chapter of psychology is usually read as eliminating normative epistemology, and Kim's objection that it changes the subject is usually taken to be decisive against it. Quine's later remarks suggest a different reading: normative epistemology survives as the technology of truth-seeking, a body of hypothetical imperatives grounded in psychology. We argue that on this reading the proposal is close to the view of EPI-26-a. It makes epistemology continuous with psychology; but its imperatives presuppose a bridge principle, that truth is the end by which methods are to be assessed, which is not a finding of psychology. Reconsidered, Quine's proposal concedes the distinctness of the normative.
\end{abstract}

\section{Two readings}

\citet{Quine1969} proposed that, since the Cartesian project of deriving knowledge of the world from sense experience had failed, epistemology should study instead how the human subject, given meagre input, produces a torrential output of theory. \citet{Kim1988} replied that this study is descriptive and that the concept of justification, which is normative, has no place in it: naturalized epistemology is not epistemology.

On the \emph{replacement reading}, Kim is right about Quine. On the \emph{engineering reading}, which Quine endorsed in replying to White, normative epistemology is retained as a branch of engineering: the technology of truth-seeking, whose norms say which methods best serve prediction \citep{Quine1986}. Its norms are hypothetical: \enquote{if you aim at true prediction, use method $M$}.

\section{The engineering reading assessed}

On the engineering reading, the antecedent of each norm is supplied by psychology and related sciences: which methods, in creatures like us and environments like ours, yield true prediction. This is Theorem~3.1 of EPI-26-a: epistemology cannot settle which norms bind human believers without psychological premises.

\begin{proposition}[The end is not a finding]\label{p:end}
The engineering reading yields categorical epistemic assessments only together with the premise that truth, or true prediction, is the end by which methods of belief formation are to be assessed. That premise is a bridge principle in the sense of Principle~P3 of EPI-26-a, and is not a finding of psychology.
\end{proposition}

\begin{proof}[Argument]
A hypothetical imperative assesses a method only relative to an end. To conclude that a subject's belief is ill-formed, and not merely ill-suited to an end she may lack, the end must be one by which belief is properly assessed. Psychology can establish which ends subjects in fact pursue and which methods serve them; by Theorem~4.1 of EPI-26-a, no set of such descriptive claims entails that one end is the proper standard of belief.
\end{proof}

Quine did not need to deny Proposition~\ref{p:end}; he took the end for granted. But taking it for granted is not deriving it, and the premise taken for granted is precisely the normative element that Kim said naturalized epistemology lacked. Reliabilists who build on psychology, such as \citet{Goldman1986}, make the same premise explicit by defining justification in terms of truth-conducive processes.

\begin{objection}
The end of truth is constitutive of belief, and so is fixed by the psychology of belief.
\end{objection}

\begin{reply}
If belief is constitutively aimed at truth, psychology may find that creatures have states with that aim. That a state ought to be assessed by its aim is a further, normative claim, and it is the claim the engineering reading needs.
\end{reply}

\section{Conclusion}

Reconsidered, Quine's proposal is not a replacement of epistemology but a demand that its norms be grounded in what psychology reveals about how belief is formed. That demand is correct, and the engineering norms it yields presuppose a normative end that psychology does not supply. Epistemology must be continuous with psychology, and its normative questions remain distinct.

\begin{thebibliography}{99}\small
\bibitem[Goldman(1986)]{Goldman1986} Goldman, A.~I. (1986). \emph{Epistemology and Cognition}. Harvard University Press.
\bibitem[Kim(1988)]{Kim1988} Kim, J. (1988). What is \enquote{naturalized epistemology}? \emph{Philosophical Perspectives}, 2, 381--405.
\bibitem[Quine(1969)]{Quine1969} Quine, W.~V. (1969). Epistemology naturalized. In \emph{Ontological Relativity and Other Essays} (pp.~69--90). Columbia University Press.
\bibitem[Quine(1986)]{Quine1986} Quine, W.~V. (1986). Reply to Morton White. In L.~E. Hahn and P.~A. Schilpp (Eds.), \emph{The Philosophy of W.~V. Quine} (pp.~663--665). Open Court.
\end{thebibliography}

\end{document}
